\subsection{Provenance}\label{subsec:provenance}

Provenance is one of the primary goals of this thesis. The system was designed not only to produce final biomedical rules, but more importantly, to make each of these final rules traceable back to the source. 

Using the back-pointers introduced in Subsection~\ref{subsec:data-flow-and-intermediate-objects}, each intermediate object can be traced back to its source through the provenance chain (Final rule $\rightarrow$ Canonical rule $\rightarrow$ Normal finding $\rightarrow$ MAP finding $\rightarrow$ Text element $\rightarrow$ Document).

This provenance chain is enforced through typed intermediate objects and back-pointers between stages. Upstream stages populate provenance fields when they produce their outputs, and downstream stages preserve these references. As a result, the output of every later stage --- normalization, grouping, and scoring --- can be followed back to the original evidence, so a reviewer can inspect a final rule, trace it back to its source text element, and compare the two.

Creating and saving typed intermediate objects enable fine-grained error analysis: if a final rule is wrong, the error can be attributed to a certain stage of the pipeline, e.g. normalization, grouping, relation classification (SUPPORT/CONTRADICT), etc.

In this thesis, we take the position that \textbf{provenance enables auditability}, but it is not a proof of correctness. A provenance chain shows where the final rule came from; it does not prove that the generated final rule is clinically valid. For instance, a final rule with a perfect provenance chain might have encoded an interpretation of a text, with which a reviewer might disagree. The distinction between auditability and correctness will be discussed in more detail in Chapter~\ref{chapter:Discussion}.

